\section{Block-angular Jacobian structure}
\label{sec:block-jacobian}

Order the variables as
\begin{equation}
  (\theta,S^{(1)},\ldots,S^{(N)})
\end{equation}
and the polynomial equations sample by sample.  Write
\begin{equation}
  C_i=\frac{\partial F^{(i)}}{\partial\theta},
  \qquad
  D_i=\frac{\partial F^{(i)}}{\partial S^{(i)}}.
\end{equation}

\begin{lemma}[Cross-sample derivatives vanish]
\label{lem:cross-derivatives}
If \(i\ne j\), then
\begin{equation}
  \frac{\partial F^{(i)}}{\partial S^{(j)}}=0.
\end{equation}
\end{lemma}

\begin{proof}
By sample locality, no indeterminate in \(S^{(j)}\) occurs in any polynomial of
\(F^{(i)}\).  A formal partial derivative with respect to an absent
indeterminate is zero.
\end{proof}

\begin{theorem}[Block-angular presentation]
\label{thm:block-jacobian}
With the order above, the global Jacobian is
\begin{equation}
J_F=
\begin{bmatrix}
C_1&D_1&0&\cdots&0\\
C_2&0&D_2&\cdots&0\\
\vdots&\vdots&\vdots&\ddots&\vdots\\
C_N&0&0&\cdots&D_N
\end{bmatrix}.
\label{eq:block-jacobian}
\end{equation}
\end{theorem}

\begin{proof}
The parameter block can be nonzero in every sample because parameters are
shared.  The local diagonal blocks contain derivatives with respect to that
sample's state.  Every off-diagonal local block is zero by
\cref{lem:cross-derivatives}.
\end{proof}

The theorem holds over the polynomial quotient ring before evaluation and over
\(\R\) after evaluation at any real point.  These are different coefficient
structures.  Gaussian elimination is immediately available for a numerical
matrix over a field.  Symbolic elimination over a multivariate quotient ring
requires additional algebra and cannot be charged as ordinary field row
reduction without justification.

\subsection{Pointwise tangent compatibility}

At a real solution \(p\), a first-order perturbation
\((\delta\theta,\delta S^{(1)},\ldots,\delta S^{(N)})\) satisfies the
linearized equations exactly when
\begin{equation}
  C_i(p)\delta\theta+D_i(p)\delta S^{(i)}=0
  \qquad (i\in[N]).
  \label{eq:linearized-local}
\end{equation}
Eliminating \(\delta S^{(i)}\) gives a constraint on
\(\delta\theta\).  This is a tangent-extension question at a known point.  It
does not ask whether a point \(p\) exists, and it does not recover the constant
terms of the original equations.

\begin{proposition}[Local linearized extension criterion]
\label{prop:local-linear-extension}
Over a field, let \(P_i^\perp\) have row space
\(\kernel D_i(p)^{\mathsf T}\).  Then
\cref{eq:linearized-local} has a solution \(\delta S^{(i)}\) if and only if
\begin{equation}
  P_i^\perp C_i(p)\delta\theta=0.
  \label{eq:local-message-criterion}
\end{equation}
\end{proposition}

\begin{proof}
The vector \(-C_i(p)\delta\theta\) belongs to the image of \(D_i(p)\) exactly
when it is annihilated by every element of the left null space of \(D_i(p)\).
Writing a basis of that null space as the rows of \(P_i^\perp\) gives
\cref{eq:local-message-criterion}.
\end{proof}

Stacking the reduced constraints yields the root matrix used by differential
CMP.  The criterion is exact for the linear system at \(p\).
